
\documentclass{beamer}

%\usepackage[table]{xcolor}
\mode<presentation> {
  \usetheme{Boadilla}

%  \usetheme{Pittsburgh}
%\usefonttheme[2]{sans}
\renewcommand{\familydefault}{cmss}
%\usepackage{lmodern}
%\usepackage[T1]{fontenc}
%\usepackage{palatino}
%\usepackage{cmbright}
  \setbeamercovered{transparent}
\useinnertheme{rectangles}
}

\setbeamercolor{normal text}{fg=black}
\setbeamercolor{structure}{fg= blue}
\definecolor{trial}{cmyk}{1,0,0, 0}
\definecolor{trial2}{cmyk}{0.00,0,1, 0}
\definecolor{darkgreen}{rgb}{0,.4, 0.1}
\usepackage{array}
\beamertemplatesolidbackgroundcolor{white}  \setbeamercolor{alerted
text}{fg=red}
\usepackage{tcolorbox}
\setbeamertemplate{caption}[numbered]\newcounter{mylastframe}


\usepackage{tikz}
\usetikzlibrary{arrows}
\usepackage{colortbl}

\renewcommand{\familydefault}{cmss}
%\usepackage[all]{xy}

\usepackage{tikz}
\usepackage{lipsum}

 \newenvironment{changemargin}[3]{%
 \begin{list}{}{%
 \setlength{\topsep}{0pt}%
 \setlength{\leftmargin}{#1}%
 \setlength{\rightmargin}{#2}%
 \setlength{\topmargin}{#3}%
 \setlength{\listparindent}{\parindent}%
 \setlength{\itemindent}{\parindent}%
 \setlength{\parsep}{\parskip}%
 }%
\item[]}{\end{list}}
\usetikzlibrary{arrows}

\usecolortheme{lily}

\newtheorem{com}{Comment}
\newtheorem{lem} {Lemma}
\newtheorem{prop}{Proposition}
\newtheorem{condition}{Condition}
\newtheorem{thm}{Theorem}
\newtheorem{defn}{Definition}
\newtheorem{cor}{Corollary}
\newtheorem{obs}{Observation}
 \numberwithin{equation}{section}
 

\setbeamertemplate{navigation symbols}{}


\title[PLSC 43401]{Mathematical Foundations of Political Methodology}

\author{Robert Gulotty}
\institute[Chicago]{University of Chicago}
\vspace{0.3in}




\begin{document}

\begin{frame}
\maketitle
\end{frame}


\begin{frame}
\frametitle{Lecture 6: Transforms}
\begin{itemize}
\item Proof that if $X_1\ ...\ X_N$ are MVN and $cov(X_i,X_j)=0$ then independent.
\item Change of Coordinates
\item Moment Generating Functions
\end{itemize}
\end{frame}





\begin{frame}

\begin{defn}
$\boldsymbol{X} = (X_{1}, X_{2}, \hdots, X_{N}) $ is a vector of random variables.  If $\boldsymbol{X}$ has pdf 

\begin{eqnarray}
f(\boldsymbol{x}) & = & (2 \pi)^{-N/2} \text{det}\left(\boldsymbol{\Sigma}\right)^{-1/2} \exp\left(-\frac{1}{2}(\boldsymbol{x} - \boldsymbol{\mu})^{'}\boldsymbol{\Sigma} (\boldsymbol{x} - \boldsymbol{\mu} ) \right) \nonumber 
\end{eqnarray}
$\boldsymbol{X}$ is a \alert{Multivariate Normal} Distribution, 
\begin{eqnarray}
\boldsymbol{X} & \sim & \text{Multivariate Normal} (\boldsymbol{\mu}, \boldsymbol{\Sigma}) \nonumber 
\end{eqnarray}


\end{defn}


\end{frame}


\begin{frame}

\begin{prop}
Suppose $\boldsymbol{X} \sim \text{Multivariate Normal}(\boldsymbol{\mu}, \boldsymbol{\Sigma})$.  
where $\boldsymbol{X}= (X_{1}, X_{2}, \hdots, X_{N})$. \\
If $\text{cov}(X_{i}, X_{j})  = 0  $, then $X_{i}$ and $X_{j}$ are independent 

\end{prop}
\end{frame}



\begin{frame}

\begin{proof}
Suppose $\mathbf{X}=(X_1,\ ... X_N)$ are MVN distributed \\ \pause 
Partition these into two groups $\mathbf{X_A}$ and $\mathbf{X_B}$, with means $\boldsymbol{\mu_A}$ and $\boldsymbol{\mu_B}$, and variance covariance 
$$\boldsymbol{\Sigma}=\ \begin{pmatrix} 
\text{var}(X_{1}) & \text{cov}(X_{1}, X_{2}) & \hdots & \text{cov}(X_{1}, X_{N}) \\
\text{cov}(X_{2}, X_{1}) & \text{var}(X_{2}) & \hdots & \text{cov}(X_{2}, X_{N} ) \\
\vdots & \vdots & \ddots & \vdots \\
\text{cov}(X_{N}, X_{1} ) & \text{cov}(X_{N}, X_{2} ) & \hdots & \text{var}(X_{N} ) \\
\end{pmatrix}=\begin{pmatrix} \boldsymbol{\Sigma}_{AA} & \boldsymbol{\Sigma}_{AB} \\ \boldsymbol{\Sigma}_{BA} &\boldsymbol{\Sigma}_{BB} \\\end{pmatrix}$$

\pause
Now assume: $\boldsymbol{\Sigma}_{AB}=\boldsymbol{\Sigma}_{BA}=0$\pause
\begin{eqnarray}
\invisible<1-3>{& & [\mathbf{x_A^T}-\boldsymbol{\mu_A^T}, \quad \mathbf{x_B^T}-\boldsymbol{\mu_B^T}]\begin{bmatrix} \boldsymbol{\Sigma}_{AA} & 0\\ 0 &\boldsymbol{\Sigma}_{BB} \end{bmatrix}^{-1}  \begin{bmatrix}\mathbf{x_A}-\boldsymbol{\mu_A} \\ \mathbf{x_B}-\boldsymbol{\mu_B}\end{bmatrix} \nonumber } 
\end{eqnarray}

\end{proof}
\end{frame}

\begin{frame}
Now taking that equation:
\begin{eqnarray}
\invisible<1>{& & [\mathbf{x_A^T}-\boldsymbol{\mu_A^T}, \quad \mathbf{x_B^T}-\boldsymbol{\mu_B^T}]\begin{bmatrix} \boldsymbol{\Sigma}_{AA} & 0\\ 0 &\boldsymbol{\Sigma}_{BB} \end{bmatrix}^{-1}  \begin{bmatrix}\mathbf{x_A}-\boldsymbol{\mu_A} \\ \mathbf{x_B}-\boldsymbol{\mu_B}\end{bmatrix} \nonumber \\ } \pause 
\invisible<1-2>{& & [\mathbf{x_A^T}-\boldsymbol{\mu_A^T}, \quad \mathbf{x_B^T}-\boldsymbol{\mu_B^T}]\begin{bmatrix} \boldsymbol{\Sigma}_{AA}^{-1} & 0\\ 0 &\boldsymbol{\Sigma}_{BB}^{-1} \end{bmatrix}  \begin{bmatrix}\mathbf{x_A}-\boldsymbol{\mu_A} \\ \mathbf{x_B}-\boldsymbol{\mu_B}\end{bmatrix} \nonumber \\ } \pause 
\invisible<1-3>{& &(\mathbf{x_A}-\boldsymbol{\mu_A})^T\boldsymbol{\Sigma}_{AA}^{-1}(\mathbf{x_A}-\boldsymbol{\mu_A})+ (\mathbf{x_B}-\boldsymbol{\mu_B})^{T} \boldsymbol{\Sigma}_{BB}^{-1}(\mathbf{x_B}-\boldsymbol{\mu_B}) \nonumber } \pause
\end{eqnarray}

\invisible<1-4>{So $exp\{(\mathbf{x}-\boldsymbol{\mu})^T\boldsymbol{\Sigma}^{-1}(\mathbf{x}-\boldsymbol{\mu})\}$ factorizes into :
$$exp\{(\mathbf{x_A}-\boldsymbol{\mu_A})^T\boldsymbol{\Sigma}_{AA}^{-1}(\mathbf{x_A}-\boldsymbol{\mu_A})\}*exp\{(\mathbf{x_B}-\boldsymbol{\mu_B})^{T} \boldsymbol{\Sigma}_{BB}^{-1}(\mathbf{x_B}-\boldsymbol{\mu_B})\}$$}
\pause
so :\pause
\invisible<1-5>{$$f(\mathbf{x})=g(\mathbf{x_A})h(\mathbf{x_B})$$}

\end{frame}


\begin{frame}
\frametitle{Change of Coordinates}


\begin{prop}
Suppose $X$ is a random variable and $Y = g(X)$, where $g:\mathbb{R} \rightarrow \mathbb{R}$ that is a $monotonic$ function.  \\
Define $g^{-1}:\mathbb{R} \rightarrow \mathbb{R}$ such that $g^{-1}(g(X)) = X$ and is differentiable.  Then, 
\begin{eqnarray}
f_{Y}(y) & = & f_{X}(g^{-1}(y))\left|\frac{\partial g^{-1}(y)}{\partial y} \right| \text{ if $y = g(x)$ for some $x$ } \nonumber \\
& = & 0 \text{ otherwise } \nonumber 
\end{eqnarray}

\end{prop}


\end{frame}

\begin{frame}
\frametitle{Change of Coordinates}


\begin{proof}

Suppose $g(\cdot)$ is monotonically increasing (WLOG)\pause 
\begin{eqnarray}
\invisible<1>{F_{Y}(y) & = & P(Y \leq y) \nonumber \\} \pause 
\invisible<1-2>{& = & P(g(X) \leq y) \nonumber \\} \pause 
\invisible<1-3>{& = & P(X \leq g^{-1}(y) ) \nonumber \\} \pause 
\invisible<1-4>{& = & F_{X}(g^{-1}(y) ) \nonumber } \pause 
\end{eqnarray}
\invisible<1-5>{Now differentiating to get the pdf } \pause 
\begin{eqnarray}
\invisible<1-6>{\frac{\partial F_{Y}(y)}{\partial y} & = & \frac{\partial F_{X}(g^{-1}(y) )} {\partial y } \nonumber \\
 & = & f_{X}(g^{-1}(y) ) \frac{\partial g^{-1}(y)}{\partial y }   \nonumber } \pause 
\end{eqnarray}
\invisible<1-7>{Then this is a pdf because $\frac{\partial g^{-1}(Y)}{\partial y } > 0$.  } 


\end{proof}



\end{frame}


\begin{frame}
\frametitle{Change of Coordinates}

Suppose $X$ is a random variable with pdf $f_{X}(x)$.  Suppose $Y =  X^{n}$.  Find $f_{Y}(y)$. \pause  \\
\invisible<1>{Then $g^{-1} (x)  =  x^{1/n}$.  } \pause 




\begin{eqnarray}
\invisible<1-2>{f_{Y}(y) & = & f_{X}(g^{-1}(y)) \left| \frac{\partial g^{-1}(Y)}{\partial y } \right| \nonumber \\} \pause 
\invisible<1-3>{ & = & f_{X} (y^{1/n} ) \frac{y^{\frac{1}{n} - 1}}{n} \nonumber } 
\end{eqnarray}
\pause
\invisible<1-4>{\alert{We've used this to derive many of the pdfs}} \pause 
\begin{itemize}
\invisible<1-5>{\item[-] Normal distribution }
\invisible<1-5>{\item[-] Chi-Squared Distribution}
\end{itemize}




\end{frame}


\begin{frame}
\frametitle{$\chi^{2}$ Distribution}

Suppose $Z \sim \text{Normal}(0,1)\leftarrow f_x(x)$. \pause   \\
\invisible<1>{Consider $X = Z^2 \leftarrow g_x(x)$} \pause 
\begin{eqnarray}
\invisible<1-2>{F_{X}(x)   & = &  P(X \leq x) \nonumber \\ } \pause 
\invisible<1-3>{& = & P(Z^2 \leq x ) \nonumber \\} \pause 
 \invisible<1-4>{& = & P(-\sqrt{x} \leq Z \leq \sqrt{x}) \leftarrow g_x^{-1}(x)\nonumber \\} \pause 
 \invisible<1-5>{& = & \frac{1}{\sqrt{2\pi}} \int_{-\sqrt{x}}^{\sqrt{x} } e^{-\frac{z^2}{2}} dz \nonumber \\} \pause 
 \invisible<1-6>{& = & F_{Z} (\sqrt{x}) - F_{Z} (-\sqrt{x}) \nonumber } 
\end{eqnarray}
\invisible<1-7>{The pdf then is } \pause 

\begin{eqnarray}
\invisible<1-8>{\frac{\partial F_{X}(x) }{\partial x }  & = & f_{Z} (\sqrt{x}) \frac{1}{2\sqrt{x}} + f_{Z}(-\sqrt{x}) \frac{1}{2\sqrt{x}} \nonumber } 
\end{eqnarray}

\pause 




\end{frame}

\begin{frame}
\frametitle{$\chi^2$ Distribution}
\begin{eqnarray}
\frac{\partial F_{X}(x) }{\partial x }  & = & f_{Z} (\sqrt{x}) \frac{1}{2\sqrt{x}} + f_{Z}(-\sqrt{x}) \frac{1}{2\sqrt{x}} \nonumber \pause \\
\invisible<1>{& = & \frac{1}{\sqrt{x}}\frac{1}{2 \sqrt{2\pi}} ( 2e^{-\frac{x}{2}}) \nonumber \\} \pause 
\invisible<1-2>{& = & \frac{1}{\sqrt{x}}\frac{1}{\sqrt{2\pi}} ( e^{-\frac{x}{2}}) \nonumber \\} \pause 
\invisible<1-3>{& = & \frac{(\frac{1}{2})^{1/2}}{\Gamma(\frac{1}{2})}\left(x^{1/2 - 1} e^{-\frac{x}{2}}\right) \nonumber } \pause 
\end{eqnarray}

\invisible<1-4>{$X \sim \text{Gamma}(1/2, 1/2)$} \pause 

\invisible<1-5>{Then if $X = \sum_{i=1}^{N} Z^2$\\}  \pause 

\invisible<1-6>{$X \sim \text{Gamma}(n/2, 1/2) $}



\end{frame}




\begin{frame}
\frametitle{Approximating functions and second order conditions}


\begin{thm}
\textbf{Taylor's Theorem}
Suppose $f:\Re \rightarrow \Re$, $f(x)$ is infinitely differentiable function.  Then, the taylor expansion of $f(x)$ around \alert{$a$} is given by 

\begin{eqnarray}
f(x) & = & f(a) + \frac{f^{'}(a)}{1!} (x- a) + \frac{f^{''}(a)}{2!} (x - a)^2 + \frac{f^{'''}(a)}{3!}(x- a)^3 + \hdots \nonumber \\
f(x) & = & \sum_{n=0}^{\infty } \frac{f^{n} (a) }{n!} (x - a)^n \nonumber
\end{eqnarray}

\end{thm}




\end{frame}


\begin{frame}
\frametitle{Example Function}

Suppose $a$ = 0 and $f(x) = e^{x}$.  Then, 
\begin{eqnarray}
f^{'}(x) & = & e^{x} \nonumber \\
f^{''}(x) & = & e^{x} \nonumber \\
\vdots  & \vdots & \vdots \nonumber \\
f^{n} (x) & = & e^{x} \nonumber 
\end{eqnarray}

This implies
\begin{eqnarray}
e^{x} & = & 1 + x + \frac{x^{2}}{2!} + \frac{x^3}{3!} \hdots + \frac{x^{n}}{n!} + \hdots \nonumber \end{eqnarray}

\end{frame}





\begin{frame}
\frametitle{Moment Generating Functions}


\begin{defn}
Suppose $X$ is a random variable with pdf $f$.  Define, 
\begin{eqnarray}
E[X^{n}] & = & \int_{-\infty}^{\infty} x^{n} f(x) dx   \nonumber 
\end{eqnarray}

We will call $X^{n}$ the \alert{$n^{\text{th}}$} moment of $X$ 

\end{defn}

\begin{itemize}
\item[-] By this definition $var(X) = \text{Second Moment} - \text{First Moment}^{2} $
\item[-] We are assuming that the integral converges
\end{itemize}

\end{frame}


\begin{frame}
\frametitle{Moment Generating Functions}

\begin{prop}
Suppose $X$ is a random variable with pdf $f(x)$.  Call $M(t) = E[e^{tX}]$, 
\begin{eqnarray}
M(t) & = & E[e^{tX}] \nonumber \\
 & = & \int_{-\infty}^{\infty} e^{tx} f(x) dx \nonumber 
 \end{eqnarray}

We will call $M(t)$ the moment generating function, because:
\begin{eqnarray}
\frac{\partial^{n} M (t) }{\partial^{n} t}|_{0} & = & E[X^{n}] \nonumber 
\end{eqnarray}

(Assuming that we can interchange derivative and integral)

\end{prop}



\end{frame}










\begin{frame}
\frametitle{Moment Generating Functions}

\begin{small}
\begin{proof}
Recall the Taylor Expansion of $e^{tX}$ at $0$, \pause 
\begin{eqnarray}
\invisible<1>{e^{tX} & = & 1 + tx + \frac{t^2 x^2}{2!} + \frac{t^3 x^3}{3!} + \hdots \nonumber } \pause 
\end{eqnarray}

\invisible<1-2>{Then, } \pause 
\begin{eqnarray}
\invisible<1-3>{E[e^{tX} ] & = &  1 + tE[X] + \frac{t^2}{2!} E[X^2] + \frac{t^3}{3!} E[X^3] + \hdots \nonumber } \pause 
\end{eqnarray}

\invisible<1-4>{Differentiate once:} \pause 
\begin{eqnarray}
\invisible<1-5>{\frac{\partial M(t)}{\partial t} & = & 0 + E[X]  + \frac{2t}{2!} E[X^2] + \hdots \nonumber \\
M^{'}(0) & = & 0 + E[X] + 0 + 0 \hdots \nonumber } 
\end{eqnarray}





\end{proof}

\end{small}



\end{frame}

\begin{frame}

\begin{small}
\begin{proof}

Differentiate $n$ times \pause 
\begin{eqnarray}
\invisible<1>{\frac{\partial^{n} M(t)}{\partial^{n} t } & = & 0 + 0 + 0 + \hdots + \frac{n \times n-1 \times \hdots 2 \times t^{0} E[X^{n}] }{n!}  + \frac{n!t E[X^{n+1}] }{(n+ 1)! } + \hdots \nonumber \\} \pause 
\invisible<1-2>{& = & \frac{n! E[X^{n}] }{n!}  + \frac{n!t E[X^{n+1}] }{(n+ 1)! } + \hdots  \nonumber } \pause 
\end{eqnarray}

\invisible<1-3>{Evaluated at 0, yields $M^{n}(0)  = E[X^{n}] $} \pause 

\end{proof}

\begin{itemize}
\invisible<1-4>{\item[-] If two random variables, $X$ and $Y$ have the same moment generating functions, then $F_{X}(x) = F_{Y}(y)$ for \alert{almost all} $x$.  } 
\end{itemize}


\end{small}
\end{frame}


\begin{frame}
\frametitle{The Moments of the Normal Distribution}
Suppose $Z \sim N(0,1)$.  \pause 

\begin{eqnarray}
\invisible<1>{E[e^{tX}] & = & \frac{1}{\sqrt{2\pi}} \int_{-\infty}^{\infty} e^{tx} e^{-x^2/2} dx \nonumber } \pause 
\end{eqnarray}

\begin{eqnarray} 
\invisible<1-2>{tx - \frac{1}{2} x^2  & = & tx - \frac{1}{2} x^2 - \frac{1}{2}t^2+\frac{1}{2}t^2 \nonumber } \\\pause 
\invisible<1-3>{ & = & -\frac{1}{2}(x-t)^2 +\frac{1}{2}t^2\nonumber   } \pause 
\end{eqnarray}
\begin{eqnarray} 
\invisible<1-4>{E[e^{tX}] & = & e^{\frac{t^{2}}{2}} \frac{1}{\sqrt{2\pi}} \int_{-\infty}^{\infty} e^{-(x- t)^2/2} dx \nonumber \\} \pause 
 \invisible<1-5>{& = & e^{\frac{t^{2}}{2}} \nonumber } 
\end{eqnarray}


\end{frame}

\begin{frame}
\frametitle{Extracting Moments of the Normal Distribution}
$$M(t)=E[e^{tX}]=e^{\frac{t^{2}}{2}}$$
\pause 
\begin{eqnarray}
\invisible<1>{M^{'}(0) & = &E[X] =  e^{t^2/2} t |_{0} = 0 \nonumber \\} \pause
\invisible<1-2>{M^{''} (0 ) & = & E[X^2] =  e^{t^2/2} (t^2 + 1) |_{0} = 1 \nonumber \\} \pause 
\invisible<1-3>{M^{'''} (0) & = & E[X^3] = e^{t^2/2} t (t^2 + 3) |_{0} = 0 \nonumber \\} \pause 
\invisible<1-4>{M^{''''} (0) & = & E[X^4] = e^{t^2/2} (t^4 + 6t^2 + 3)|_{0} = 3 \nonumber } \\ \pause 
\invisible<1-5>{M^{5} (0) & = & E[X^5] = e^{t^2/2} t (t^4 + 10t^2 + 15)|_{0} = 0 \nonumber } \\ \pause 
\invisible<1-6>{M^{6} (0) & = & E[X^6] = e^{t^2/2} (t^6 + 15t^4 + 45t^2 + 15 )|_{0} = 15 \nonumber } \pause 
\end{eqnarray} 


\end{frame}

\begin{frame}

\begin{prop}
Suppose $X_{i}$ are a sequence of independent random variables.  Define 
\begin{eqnarray}
Y & = & \sum_{i=1}^{N} X_{i} \nonumber 
\end{eqnarray}

Then 
\begin{eqnarray}
M_{Y}(t) & = & \prod_{i=1}^{N} M_{X_{i}}(t) \nonumber 
\end{eqnarray}


\end{prop}

\end{frame}

\begin{frame}

\begin{proof}
\begin{eqnarray}
M_{Y}(t) & = & E[e^{tY}] \nonumber \\ \pause 
\invisible<1>{& = & E[e^{t\sum_{i=1}^{N} X_{i} } ] \nonumber \\ } \pause 
 \invisible<1-2>{& = & E[e^{tX_{1} + tX_{2} + \hdots tX_{N} } ] \nonumber \\} \pause 
  \invisible<1-3>{& = & E[e^{tX_{1} }]E[e^{tX_{2} }] \hdots E[e^{tX_{N} }] \text{ (by independence) } \nonumber \\} \pause 
   \invisible<1-4>{& = & \prod_{i=1}^{N} E[e^{tX_{i}}] \nonumber } 
\end{eqnarray}

\end{proof}

\end{frame}

\begin{frame}
\frametitle{Inequalities and Limit Theorems}


\alert{Limit Theorems}
\begin{itemize}
\item[-] What happens when we consider a long sequence of random variables?
\item[-] What can we reasonably infer from data?
\begin{itemize}
\item[-] Laws of large numbers: averages of random variables converge on expected value?
\item[-] Central Limit Theorems: sum of random variables have normal distribution?
\end{itemize}
\item[-] We'll focus on intuition for both, but we'll prove some stuff too.  
\end{itemize}

\alert{Review Session}


\end{frame}


\begin{frame}
\frametitle{Weak Law of Large Numbers}

Proof plan: 
\begin{itemize}
\item[-] Markov's Inequality
\item[-] Chebyshev's Inequality
\item[-] Weak Law of Large Numbers 
\end{itemize}


\end{frame}



\begin{frame}
\frametitle{Markov's Inequality}

\begin{prop}
Suppose $X$ is a random variable that takes on non-negative values.  Then, for all $a>0$, 
\begin{eqnarray}
P(X\geq a) & \leq & \frac{E[X]}{a} \nonumber 
\end{eqnarray}
\end{prop}

\end{frame}


\begin{frame}
\frametitle{Markov's Inequality} 


\begin{proof}
For $a>0$,  \pause 
\begin{eqnarray}
\invisible<1>{E[X] & = & \int_{0}^{\infty} x f(x) dx \nonumber \\} \pause 
 \invisible<1-2>{& = & \int_{0}^{a} x f(x) dx + \int_{a}^{\infty} x f(x) dx \nonumber } \pause 
\end{eqnarray}

\invisible<1-3>{Because $X\geq 0$, } \pause 

\begin{eqnarray}
\invisible<1-4>{E[X] \geq \int_{a}^{\infty} x f(x) dx  \geq \int_{a}^{\infty} a f(x)dx = a P(X \geq a )\nonumber \\} \pause 
\invisible<1-5>{\frac{E[X]}{a} \geq P(X \geq a )\nonumber } 
\end{eqnarray}


\end{proof}



\end{frame}

\begin{frame}
\frametitle{Chebyshev's Inequality} 

\begin{prop} 
If $X$ is a random variable with mean $\mu$ and variance $\sigma^2$, then, for any value $k>0$, 
\begin{eqnarray}
P(|X - \mu| \geq k) & \leq & \frac{\sigma^2}{k^2} \nonumber 
\end{eqnarray}

\end{prop}


\end{frame}

\begin{frame}
\frametitle{Chebyshev's Inequality} 


\begin{proof}
Define the random variable 
\begin{eqnarray}
 Y & = & (X - \mu)^2 \nonumber 
\end{eqnarray}

Where $\mu = E[X]$. 

\pause 

\invisible<1>{Then we know $Y$ is a non-negative random variable.  Set $a = k^2$.  \\} \pause 
\invisible<1-2>{Applying the inequality:} \pause 
\begin{eqnarray}
\invisible<1-3>{P(Y \geq k^2) \leq \frac{E[Y]}{k^2} \nonumber \\} \pause 
\invisible<1-4>{P( (X- \mu)^2 \geq k^2) \leq \frac{E[(X - \mu)^2]}{k^2} \nonumber \\} \pause 
\invisible<1-5>{P( (X- \mu)^2 \geq k^2) \leq \frac{\sigma^2}{k^2}\nonumber } 
\end{eqnarray}

\end{proof}

\end{frame}


\begin{frame}
\frametitle{Chebyshev's Inequality} 


Further we know that, 
\begin{eqnarray}
(X - \mu)^2  \geq k^2 \nonumber 
\end{eqnarray}

\pause 

\invisible<1>{Implies that } \pause 
\begin{eqnarray}
\invisible<1-2>{|X - \mu| \geq k \nonumber } \pause 
\end{eqnarray}

\invisible<1-3>{Thus, we have shown } \pause 
\begin{eqnarray}
\invisible<1-4>{P( |X- \mu| \geq k) \leq \frac{\sigma^2}{k^2}\nonumber } 
\end{eqnarray}

\end{frame}



\begin{frame}
\frametitle{Sequence of Random Variables}

Sequence of Independent and Identically, Distributed Random variables.
\begin{itemize}
\item[-] Sequence: $X_{1}, X_{2}, \hdots, X_{n}, \hdots $
\item[-] Think of a sequence as sampled \alert{data}:
\begin{itemize}
\item[-] Suppose we are drawing a sample of $N$ observations
\item[-] Each observation will be a random variable, say $X_{i}$
\item[-] With realization $x_{i}$ 
\end{itemize}
\end{itemize}

\end{frame}


\begin{frame}
\frametitle{Mean/Variance of Sample Mean}

\begin{prop}
Let $X_{1}, X_{2}, \hdots, X_{n}$ be a random sample from a distribution with mean $\mu$ and variance $\sigma^2$.  Let $\bar{X}_{n}$ be the sample mean.  Then 
$E[\bar{X}_{n}] = \mu$ and var$(\bar{X}_{n}) = \frac{\sigma^2}{n}$
\end{prop}

\pause 
\begin{proof}

\begin{eqnarray}
\invisible<1>{E[\bar{X}_{n}] & = & \frac{1}{n}\sum_{i=1}^{n} E[X_{i}] \nonumber \\} \pause 
\invisible<1-2>{& = & \frac{1}{n} n \mu  = \mu \nonumber } 
\end{eqnarray}



\end{proof}


\end{frame}


\begin{frame}
\frametitle{Mean/Variance of Sample Mean}


\pause 
\begin{eqnarray}
\invisible<1>{\text{var}(\bar{X}_{n}) & = & \frac{1}{n^2} \text{var}(\sum_{i=1}^{n} X_{i}) \nonumber \\} \pause 
 \invisible<1-2>{&= & \frac{1}{n^2} \sum_{i=1}^{n} \text{var}(X_{i}) \nonumber \\} \pause 
  \invisible<1-3>{& = & \frac{1}{n^2} n \sigma^2 = \frac{\sigma^2}{n} \nonumber } 
\end{eqnarray}


\end{frame}



\begin{frame}
\frametitle{Weak Law of Large Numbers} 

\begin{prop}
Suppose $X_{1}, X_{2}, \hdots, X_{n}$ is a random sample from a distribution with mean $\mu$ and $Var(X_{i})= \sigma^2$.  Then, for all $\epsilon >0$, 
\begin{eqnarray}
P\left\{ \left| \frac{X_{1} + X_{2} + \hdots + X_{n} }{n} -\mu \right| \geq \epsilon \right\} \rightarrow 0 \text{ as } n \rightarrow \infty \nonumber 
\end{eqnarray}

\end{prop}


\end{frame}


\begin{frame}
\frametitle{Weak Law of Large Numbers}

\begin{proof}
From our previous proposition \pause 
\begin{eqnarray}
\invisible<1>{\frac{\text{E}[X_{1} + X_{2} + \cdots + X_{n} ]}{n} & = & \frac{\sum_{i=1}^{n} E[X_{i}] }{n}  = \mu \nonumber} \pause 
\end{eqnarray}
\invisible<1-2>{Further, } \pause 
\begin{eqnarray}
\invisible<1-3>{\text{E}[ (\frac{\sum_{i=1}^{n} X_{i} - \mu}{n} )^2] = \frac{\text{Var}(X_{1} + X_{2} + \cdots + X_{n} )}{n^2} & = & \frac{ \sum_{i=1}^{n} \text{Var}(X_{i}) }{n^2} = \frac{\sigma^2}{n} \nonumber } \pause 
\end{eqnarray}
\invisible<1-4>{Apply Chebyshev's Inequality: } \pause 
\begin{eqnarray}
\invisible<1-5>{P\left\{ \left| \frac{X_{1} + X_{2} + \hdots + X_{n} }{n} -\mu \right| \geq \epsilon \right\}\leq \frac{\sigma^2}{n \epsilon^2} \nonumber }
\end{eqnarray}

\end{proof}


\end{frame}

\begin{frame}
Suppose $X_{1}, X_{2}, \hdots$ are iid normal distributions,
\begin{eqnarray}
X_{i} & \sim & \text{Normal}(0, 10) \nonumber 
\end{eqnarray}

\begin{eqnarray} 
P\left\{ \left| \frac{X_{1} + X_{2} + \hdots + X_{n} }{n} -\mu \right| \geq 0.1 \right\} \text{ as } n \rightarrow \infty\nonumber 
\end{eqnarray}
\pause 


\only<3-6>{Suppose we want to guarantee that we have at most a $0.01$ probability of being more than $0.1$ away from the true $\mu$.  How big do we need $n$?\pause 
\begin{eqnarray}
\invisible<1-3>{0.01 & = & \frac{10}{n (0.1^2) } \nonumber} \pause  \\
\invisible<1-4>{n & = & \frac{1000}{0.01} \nonumber } \pause \\
\invisible<1-5>{n & = & 100,000 } \nonumber 
\end{eqnarray}
}


\end{frame}

\end{document}